%%% Introducción %%%
\section{Introducción}

La pobreza es un fenómeno social que se manifiesta como la falta de recursos para cubrir necesidades básicas y está vinculada con factores estructurales, culturales y políticos que afectan la cohesión social, el crecimiento económico y la dignidad humana. Como plantean \cite{sen1999} y \cite{townsend1979}, va más allá de la carencia de ingresos: implica una privación de capacidades y oportunidades. Por ello, los Estados deben abordar sus causas estructurales mediante políticas públicas integrales. Sin embargo, la focalización es quizá uno de los retos más importantes, dado que requiere contar con información precisa y actualizada, así como con mecanismos institucionales capaces de distinguir correctamente a quienes realmente se encuentran en situación de pobreza, evitando tanto la exclusión de beneficiarios legítimos como la inclusión de quienes no lo necesitan. Dado que las herramientas tradicionales, como encuestas y censos, resultan costosas y de baja frecuencia, el uso de modelos de aprendizaje automático (ML) ofrece una alternativa prometedora al aprovechar su capacidad para predecir clases con alta precisión a partir de grandes volúmenes de datos. Este trabajo, por tanto, se propone desarrollar un modelo de predicción de hogares en condición de pobreza, con el fin de contribuir a soluciones más eficientes para su identificación y caracterización.

En 2024, el Departamento Administrativo Nacional de Estadística - de ahora en adelante DANE - reportó que la pobreza monetaria en Colombia fue de 31,8\% y la pobreza extrema de 11,7\% \cite{DANE2025PobrezaMonetaria}. Aunque estos niveles siguen siendo elevados, los indicadores muestran una mejora gradual. Sin embargo, las cifras nacionales ocultan marcadas disparidades regionales: departamentos como Chocó, La Guajira y Sucre registran tasas que duplican el promedio del país. Además, la pobreza está estancada en las zonas rurales y afecta con mayor intensidad a comunidades indígenas, afrocolombianas y migrantes venezolanos \cite{WorldBank2022ColombiaPoverty}. Estas diferencias evidencian una heterogeneidad profunda en el territorio nacional y señalan la necesidad de una ruta clara de política pública que priorice la atención a las poblaciones más vulnerables.

Uno de los enfoques más comunes para medir la pobreza es el basado en la falta de ingresos, aunque presenta limitaciones al depender de encuestas y censos costosos y poco frecuentes. Según \cite{SosaEscuderoAnauatiBrau2022} y \cite{Jean2016}, los modelos de aprendizaje automático pueden superar estas restricciones al integrar diversas fuentes de datos como encuestas, registros administrativos e imágenes satelitales para mejorar la precisión de las estimaciones. 

%resumen de la competencia en kaggle
Para este ejercicio de predicción\footnote{\href{https://github.com/alejandro-cardiles/PSET_2_BDML_2025_02_G2}{Enlace de replicabilidad del ejercicio}}, se tomaron los datos del DANE de la Misión de Empalme de las Series de Empleo, Pobreza y Desigualdad (MESE) para el año 2018, idóneos porque permiten observar al hogar en diferentes dimensiones como ocupación laboral, edades, actividades y geografía. En este marco, se implementaron diversos algoritmos de clasificación supervisada, entre ellos Regresión Logística (Logit) \cite{cox1958logit}, Elastic Net \cite{zou-2005}, CART \cite{breiman1984cart}, Random Forest \cite{breiman-2001} y XGBoost \cite{chen2016xgboost}, con el propósito de predecir la condición de pobreza monetaria a partir de características socioeconómicas de los hogares. La evaluación del desempeño se realizó mediante la optimización del F1, una métrica especialmente adecuada para contextos con clases desbalanceadas, al equilibrar la precisión (precision) y la sensitividad (recall). Entre los modelos estimados, XGBoost presentó el mejor rendimiento, alcanzando un F1-score promedio del \textbf{71,6\%} , superando consistentemente a los modelos lineales y de árboles simples. Este resultado se atribuye a las ventajas estructurales del algoritmo, que incluyen la capacidad de capturar relaciones no lineales y efectos de interacción complejos, la ponderación iterativa de errores residuales y la optimización del gradiente para mejorar la convergencia.

Estos resultados muestran que modelos no lineales como XGBoost pueden contribuir de manera efectiva a mejorar la focalización de las políticas sociales, especialmente en contextos caracterizados por una alta heterogeneidad y un marcado desbalance de clases.

%conclusiones finales

